Uma forma simples e intuitiva de aplicar o pareamento é o chamado \textbf{pareamento exato}. 
A lógica é direta: para cada participante do programa, buscamos um não participante que seja idêntico em todas as características relevantes observáveis — por exemplo, mesma idade, mesmo nível educacional, mesmo gênero e mesma região de residência. 
Ao comparar esses pares de indivíduos idênticos, a única diferença entre eles seria o fato de um ter recebido o tratamento (curso de qualificação) e o outro não.  

No exemplo do programa de qualificação, suponha que tenhamos dois jovens com exatamente o mesmo perfil socioeconômico, ambos com 18 anos, ensino médio incompleto, morando no mesmo bairro. 
Um deles se inscreveu no curso, o outro não. 
Se, após o curso, o primeiro conseguiu emprego e o segundo não, essa diferença pode ser atribuída, em grande parte, ao impacto do programa, já que os fatores observáveis eram iguais.  

A intuição pode ser resumida com uma tabela simples:

\begin{center}
\begin{tabular}{lccccc}
\hline
Indivíduo & Idade & Escolaridade & Bairro & Participação no Curso & Emprego\\
\hline
A (tratado) & 18 & Médio incompleto & Bairro X & Sim & Sim \\
B (controle) & 18 & Médio incompleto & Bairro X & Não & Não \\
\hline
\end{tabular}
\end{center}

Nesse exemplo, o indivíduo A e o indivíduo B são comparáveis em todas as variáveis listadas. 
A única diferença é a participação no programa, o que nos permite interpretar a discrepância no resultado (emprego após 6 meses) como efeito do curso.  

\subsection*{Vantagens e Limitações}
\textbf{Vantagem:} o pareamento exato é altamente intuitivo, transparente e fácil de explicar a públicos não técnicos.  
\textbf{Limitação:} na prática, é raro encontrar pares exatamente idênticos quando usamos muitas variáveis ao mesmo tempo. 
Esse problema é conhecido como a \textit{maldição da dimensionalidade}: quanto mais características relevantes incluímos, mais difícil se torna encontrar comparações exatas.  

Por essa razão, embora o pareamento exato seja conceitualmente poderoso e útil em casos simples, avaliações reais de políticas públicas geralmente recorrem a variações mais flexíveis, conhecidas como \textbf{métodos de pareamento aproximado}, que discutiremos na próxima seção.

\color{red}

Seria interessante fazer uma imagem para melhorar a intuição. Fazer 8 quadrados do lado esquerdo, representando os indivíduos tratados, e 16 do lado direito, representando as possíveis pessoas no grupo de comparação. Pintar de duas cores diferentes para representar homem/mulher e colocar a idade dentro do quadrado. Fazer desbalanceado, tal que a comparação das médias geraria viés de seleção. Fazer uma segunda imagem fazendo o pareamento exato nesse caso.

\color{black}